\par\noindent\textbf{(1)} 各区域能量相对势为 ε\_j=E−V\_j，Dirac色散 ε\_j²=m²c⁴+c²p\_j²。对于实p，平面波自旋or满足 (cpσ\_x+mc²σ\_z)u=εu，概率流 j=cψ†σ\_xψ；必须选群速度方向与能量分支相匹配的出射解。
\par\noindent\textbf{(2)} 取自旋or u(ε,p)∝(1,cp/(ε+mc²))\textasciicircum{}T，在x=0连续 ψ\_in+ψ\_ref=ψ\_trans。设入射上分量1，反射上分量r，透射上分量t=1+r；若η\_L=cp\_L/(E+mc²)，η\_R=cp\_R/(E−V₀+mc²)（p\_R要按向右群速度选择），第二分量连续给 η\_L(1−r)=η\_R(1+r)，故 r=(η\_L−η\_R)/(η\_L+η\_R)。
\par\noindent\textbf{(3)} 若|E−V₀|<mc²，右区p\_R纯虚，为指数衰减禁阻区，无向右远场传播；若E−V₀<−mc²，即V₀>E+mc²，右区进入负能连续谱但可以传播，这就是Klein区。对负能支，群速度 ∂ε/∂p=c²p/ε，向右传播需p\_R<0，而不是机械取正根。
\par\noindent\textbf{(4)} 单粒子Dirac概率流仍由连续方程守恒，可写 j\_in+j\_ref=j\_trans（反射j\_ref<0）；不同出射负能支的系数若用非匹配流归一化，会呈反射大于1的Klein悖论。强超临界外场的真正物理解读涉及电子—正电子对产生和多体真空，单粒子步骤只提供散射通道结构。
